\textbf{Phần 1. Trắc nghiệm lựa chọn}
\setcounterpageref{ex}{0}

\begin{ex} %[2D4V24]
Tính giá trị của biểu thức: $48 + 37 - 29$.
\choice
{$54$}
{\True $56$}
{$55$}
{$57$}
\loigiai{
$48 + 37 = 85$; $85 - 29 = 56$.
\textbf{Đáp án của câu hỏi là B}
}
\end{ex}

\begin{ex} %[2D4V24]
Một đội xe có $54$ xe tải và xe con, trong đó có $28$ xe tải. Hỏi đội xe có bao nhiêu xe con?
\choice
{\True $26$ xe con}
{$24$ xe con}
{$25$ xe con}
{$28$ xe con}
\loigiai{
$54 - 28 = 26$ (xe con).
\textbf{Đáp án của câu hỏi là A}
}
\end{ex}

\begin{ex} %[2D4V24]
Tìm $x$ biết: $x - 39 = 62 - 18$.
\choice
{$81$}
{$82$}
{\True $83$}
{$84$}
\loigiai{
Vế phải: $62 - 18 = 44$. Ta có: $x = 44 + 39 = 83$.
\textbf{Đáp án của câu hỏi là C}
}
\end{ex}

\begin{ex} %[2D4V24]
Đường gấp khúc $ABCD$ gồm $3$ đoạn thẳng: $AB = 25\text{ cm}$, $BC = 38\text{ cm}$, $CD = 29\text{ cm}$. Chiều dài đường gấp khúc $ABCD$ là
\choice
{$90\text{ cm}$}
{$91\text{ cm}$}
{$93\text{ cm}$}
{\True $92\text{ cm}$}
\loigiai{
$25 + 38 + 29 = 92\text{ (cm)}$.
\textbf{Đáp án của câu hỏi là D}
}
\end{ex}

\begin{ex} %[2D4V24]
Một cửa hàng có $92\text{ L}$ dầu, ngày thứ nhất bán được $35\text{ L}$, ngày thứ hai bán được $28\text{ L}$. Hỏi cửa hàng còn lại bao nhiêu lít dầu?
\choice
{$28\text{ L}$}
{\True $29\text{ L}$}
{$30\text{ L}$}
{$27\text{ L}$}
\loigiai{
Tổng bán: $35 + 28 = 63\text{ (L)}$. Còn lại: $92 - 63 = 29\text{ (L)}$.
\textbf{Đáp án của câu hỏi là B}
}
\end{ex}

\begin{ex} %[2D4V24]
Tìm số tự nhiên lớn nhất có hai chữ số mà khi trừ đi $36$ ta được kết quả nhỏ hơn $30$.
\choice
{\True $65$}
{$64$}
{$66$}
{$63$}
\loigiai{
Ta có: $x - 36 < 30 \Rightarrow x < 66$. Số tự nhiên lớn nhất là 65.
\textbf{Đáp án của câu hỏi là A}
}
\end{ex}

\begin{ex} %[2D4V24]
Bác An nuôi $62$ con gà và vịt, trong đó có $35$ con gà. Hỏi số vịt ít hơn số gà bao nhiêu con?
\choice
{$7$ con}
{$9$ con}
{\True $8$ con}
{$6$ con}
\loigiai{
Số vịt là: $62 - 35 = 27$ (con).\nSố vịt ít hơn số gà: $35 - 27 = 8$ (con).
\textbf{Đáp án của câu hỏi là C}
}
\end{ex}

\begin{ex} %[2D4V24]
Điền số thích hợp: $73 - \dots = 28 + 19$.
\choice
{$24$}
{$25$}
{$27$}
{\True $26$}
\loigiai{
Vế phải: $28 + 19 = 47$. Ta có: $73 - 47 = 26$.
\textbf{Đáp án của câu hỏi là D}
}
\end{ex}

\begin{ex} %[2D4V24]
Chu vi của một hình tứ giác có độ dài các cạnh là $18\text{ cm}, 24\text{ cm}, 19\text{ cm}, 25\text{ cm}$ là
\choice
{$84\text{ cm}$}
{\True $86\text{ cm}$}
{$85\text{ cm}$}
{$88\text{ cm}$}
\loigiai{
$18 + 24 + 19 + 25 = 86\text{ (cm)}$.
\textbf{Đáp án của câu hỏi là B}
}
\end{ex}

\begin{ex} %[2D4V24]
Tìm một số biết rằng lấy số đó cộng với $37$ rồi trừ đi $25$ thì được $68$.
\choice
{\True $56$}
{$54$}
{$55$}
{$58$}
\loigiai{
Tính ngược từ cuối: $68 + 25 - 37 = 93 - 37 = 56$.
\textbf{Đáp án của câu hỏi là A}
}
\end{ex}

\begin{ex} %[2D4V24]
Hai bao gạo nặng tất cả $80\text{ kg}$. Nếu chuyển $5\text{ kg}$ từ bao thứ nhất sang bao thứ hai thì hai bao nặng bằng nhau. Hỏi ban đầu bao thứ nhất nặng bao nhiêu ki-lô-gam?
\choice
{$40\text{ kg}$}
{$50\text{ kg}$}
{\True $45\text{ kg}$}
{$42\text{ kg}$}
\loigiai{
Khi bằng nhau mỗi bao nặng: $80 : 2 = 40\text{ (kg)}$.\nBan đầu bao thứ nhất nặng: $40 + 5 = 45\text{ (kg)}$.
\textbf{Đáp án của câu hỏi là C}
}
\end{ex}

\begin{ex} %[2D4V24]
Điền dấu thích hợp: $81 - 46 \dots 90 - 55$.
\choice
{$>$}
{$<$}
{Không so sánh được}
{\True $=$}
\loigiai{
$81 - 46 = 35$; $90 - 55 = 35$. Hai vế bằng nhau nên điền dấu $=$.
\textbf{Đáp án của câu hỏi là D}
}
\end{ex}

\textbf{Phần 2. Trắc nghiệm đúng sai}
\setcounterpageref{ex}{0}

\begin{ex} %[2D4V24]
Xét tính đúng sai của các bài toán rèn luyện tư duy Chủ đề 4:
\choiceTF
{\True Tổng của $49 + 36$ bằng hiệu của $95 - 10$.}
{Biểu thức $82 - 38 + 16$ có kết quả là 50.}
{\True Trong phép trừ có nhớ trong phạm vi 100, số trừ có thể lớn hơn hoặc bằng 10.}
{\True Khi cùng thêm vào số bị trừ và số trừ một số đơn vị thì hiệu vẫn giữ nguyên.}
\loigiai{
\begin{itemize}
    \item Ý a) Đúng: $49 + 36 = 85$; $95 - 10 = 85$. Bằng nhau.
    \item Ý b) Sai: $82 - 38 = 44$; $44 + 16 = 60$ (không phải 50).
    \item Ý c) Đúng: Đúng.
    \item Ý d) Đúng: Tính chất của phép trừ.
\end{itemize}
\begin{itemize}
    \item a) \textbf{Đúng}.
    \item b) \textbf{Sai}.
    \item c) \textbf{Đúng}.
    \item d) \textbf{Đúng}.
\end{itemize}
}
\end{ex}

\begin{ex} %[2D4V24]
Một cửa hàng có 80 cuộn len gồm ba màu đỏ, vàng, xanh. Biết có 28 cuộn màu đỏ và 35 cuộn màu vàng:
\choiceTF
{\True Tổng số cuộn len đỏ và vàng là 63 cuộn.}
{\True Số cuộn len màu xanh là 17 cuộn.}
{Số cuộn len màu xanh nhiều hơn số cuộn len màu đỏ.}
{Cửa hàng có số cuộn len màu vàng gấp đôi số cuộn len màu đỏ.}
\loigiai{
\begin{itemize}
    \item Ý a) Đúng: $28 + 35 = 63$ cuộn.
    \item Ý b) Đúng: $80 - 63 = 17$ cuộn.
    \item Ý c) Sai: Xanh có 17 cuộn, đỏ có 28 cuộn nên xanh ít hơn đỏ.
    \item Ý d) Sai: 35 không gấp đôi 28. Khẳng định sai.
\end{itemize}
\begin{itemize}
    \item a) \textbf{Đúng}.
    \item b) \textbf{Đúng}.
    \item c) \textbf{Sai}.
    \item d) \textbf{Sai}.
\end{itemize}
}
\end{ex}

\textbf{Phần 3. Trả lời ngắn}
\setcounterpageref{ex}{0}

\begin{ex} %[2D4V24]
Tính giá trị biểu thức: $73 - 38 + 19$.
\shortans{54}
\loigiai{
\textbf{Đáp số:} $54$\\
$73 - 38 = 35$; $35 + 19 = 54$.
}
\end{ex}

\begin{ex} %[2D4V24]
Mẹ mua 45 quả trứng, đã dùng hết 18 quả. Số trứng còn lại mẹ chia đều cho 3 người con. Hỏi mỗi người con được mấy quả trứng?
\shortans{9}
\loigiai{
\textbf{Đáp số:} $9$\\
Số trứng còn lại: $45 - 18 = 27$ quả. Mỗi người được: $27 : 3 = 9$ quả.
}
\end{ex}

\begin{ex} %[2D4V24]
Tìm số tự nhiên x biết: $84 - x = 27 + 19$.
\shortans{38}
\loigiai{
\textbf{Đáp số:} $38$\\
$27 + 19 = 46 \Rightarrow x = 84 - 46 = 38$.
}
\end{ex}

\begin{ex} %[2D4V24]
Tính chu vi hình tứ giác có độ dài các cạnh là 22 cm, 28 cm, 17 cm, 23 cm.
\shortans{90}
\loigiai{
\textbf{Đáp số:} $90$\\
Chu vi: $(22 + 28) + (17 + 23) = 50 + 40 = 90\text{ cm}$.
}
\end{ex}

\begin{ex} %[2D4V24]
Tìm một số biết rằng lấy 92 trừ đi số đó rồi cộng với 15 thì được 49.
\shortans{58}
\loigiai{
\textbf{Đáp số:} $58$\\
$92 - x = 49 - 15 = 34 \Rightarrow x = 92 - 34 = 58$.
}
\end{ex}

\begin{ex} %[2D4V24]
Hai thùng có tất cả 75 lít dầu. Thùng thứ nhất hơn thùng thứ hai 9 lít dầu. Hỏi thùng thứ nhất có bao nhiêu lít dầu?
\shortans{42}
\loigiai{
\textbf{Đáp số:} $42$\\
Thùng 1 có: $(75 + 9) : 2 = 42\text{ L}$.
}
\end{ex}
